(1) 将给定方程整理为 \(p=RT/(v-b)-a/v^2\)，可容许的体积必须 \(v>b\)。临界等温线上液体与气体分支合并为水平拐点，因此同时满足 \((\partial p/\partial v)_{T_c}=0\)、\((\partial^2p/\partial v^2)_{T_c}=0\)；前者为 \(-RT/(v-b)^2+2a/v^3\)，后者为 \(2RT/(v-b)^3-6a/v^4\)。
(2) 将两导数为零相除消去温度，得到 \(2/(v_c-b)=3/v_c\)，即 \(v_c=3b\)。再代回一阶条件得 \(RT_c=8a/(27b)\)，故 \(T_c=8a/(27Rb)\)。带回状态方程得 \(p_c=a/(27b^2)\)，并有 \(Z_c=p_cv_c/(RT_c)=3/8\)，结果与 \(a,b\) 数值无关但只适用于这一模型。
(3) 置 \(p=p_cp_r\)、\(v=3bv_r\)、\(T=T_cT_r\)，代入并逐项约去 \(a,b\)，得到 \(p_r=8T_r/(3v_r-1)-3/v_r^2\)。等价写成 \((p_r+3/v_r^2)(3v_r-1)=8T_r\)。这个普适约化形式只是在范德瓦耳斯模型内的对应态关系，不表示所有真实气体恰遵循 \(Z_c=3/8\)。
(4) 自旋odal由 \((\partial p/\partial v)_T=0\) 给定，在 \(T<T_c\) 时出现两个边界根，中间区域 \((\partial p/\partial v)_T>0\)。因等温压缩率 \(\kappa_T=-v^{-1}(\partial v/\partial p)_T\)，此段为负，与稳定系统对压缩扰动的正响应矛盾。真正液汽共存压强由化学势相等或 Maxwell 等面积构造决定，一般位于两个自旋odal压强之间而不是其中任一边界。
